\uuid{sZGa}
\exo7id{2510}
\titre{exo7 2510}
\auteur{mayer}
\organisation{exo7}
\datecreate{2009-04-01}
\isIndication{false}
\isCorrection{false}
\chapitre{Différentiabilité, calcul de différentielles}
\sousChapitre{Différentiabilité, calcul de différentielles}
\cpge{mpsi-mp2i.fonctions-deux-variables.a, pcsi.fonctions-deux-variables.a, ptsi.fonctions-deux-variables.a, bcpst1.fonctions-deux-variables.a, mp-mpi.calcul-differentiel-optimisation.b, pc.fonctions-plusieurs-variables.c, psi.fonctions-plusieurs-variables.c}
\module{Calcul différentiel}
\niveau{L3}
\difficulte{}

\contenu{
\question{
Soit $\alpha >0$. \'Etudier la
diff\'erientiabilit\'e \`a l'origine de l'application $f:\Rr^2 \to
\Rr$ qui est d\'efinie par $f(0,0) =0$ et par
$$ f(x,y) = \frac{|xy|^\alpha}{\sqrt{x^2+3y^2}} \quad \text{si} \;\;
(x,y) \neq (0,0)\; .$$
}
}
